\chapter{$\Spec(R)$: from a ring to a space}\label{ch:ch01-spec}

\noindent\textbf{Win:} turn any ring into a topological space and read off its
points.\quad\textbf{Time:} \textasciitilde15 min.\quad\textbf{Needs:} prime
ideals, radicals, localization.

\medskip

You already know the algebra of a commutative ring $R$ --- primes, radicals,
localization. This lesson installs the one idea that converts \emph{all} of that
knowledge into geometry: the \term{spectrum} $\Spec(R)$. Everything in scheme
theory is built on top of this single move, so it's the right first
brick.\scite{00DZ}{Primary reference: Stacks Project, \emph{Algebra}, ``The spectrum of a ring.''}

\section{The points are the primes}

\begin{defbox}[Definition --- the spectrum]
$\Spec(R)$ is the \textbf{set of prime ideals} of $R$. Each prime $\fp$ is a
\term{point} of the space.\scite{00DZ}{Recall a prime $\fp\neq R$ satisfies $ab\in\fp\Rightarrow a\in\fp$ or $b\in\fp$.}
\end{defbox}

Why primes and not, say, maximal ideals? Because we want the construction to be
\textbf{functorial}: a ring map $\varphi\colon R\to S$ should give a map of
spaces. It does --- $\fq\mapsto \varphi^{-1}(\fq)$ --- precisely because the
preimage of a prime is always prime. (Preimages of maximal ideals need not be
maximal.)\scite{00E2}{Functoriality of $\Spec$.} Note the arrow \emph{reverses}:
$\Spec$ is contravariant.
\[
\begin{tikzcd}[column sep=large]
R \arrow[r, "\varphi"] & S \\
\Spec(R) & \Spec(S) \arrow[l, "\varphi^{-1}"']
\end{tikzcd}
\]

\begin{intubox}[Intuition --- a ``point'' is a way of evaluating]
Think of elements of $R$ as \emph{functions} on the space. At a point $\fp$, the
value of $f$ lives in the residue field $\res(\fp)=\Frac(R/\fp)$, and
``$f(\fp)=0$'' means exactly $f\in\fp$. A prime ideal is the set of functions
vanishing at that point --- that's why primes deserve to be called points.
\end{intubox}

\section{The closed sets are $V(I)$}

\begin{defbox}[Definition --- vanishing sets]
For an ideal $I\subseteq R$, $\;V(I)=\{\fp\in\Spec(R)\mid I\subseteq\fp\}$ ---
the points where ``every function in $I$ vanishes.'' The \textbf{Zariski
topology} declares the sets $V(I)$ to be exactly the closed
sets.\scite{00E0}{These really form a topology: arbitrary $\cap$, finite $\cup$, $\emptyset$ and the whole space.}
\end{defbox}

Two facts make this a faithful dictionary back to algebra:
\begin{itemize}[leftmargin=1.4em,itemsep=2pt]
  \item \textbf{Only the radical matters:} $V(I)=V(\sqrt I)$, since
  $\fp\supseteq I \iff \fp\supseteq\sqrt I$. So radical ideals
  $\leftrightarrow$ closed sets, order-\emph{reversingly}.\scite{00E5}{$V(I)=V(\sqrt I)$ and the closed-set calculus.}
  \item \textbf{Inclusion reverses:} $I\subseteq J \Rightarrow V(I)\supseteq
  V(J)$. More constraints $=$ fewer points.
\end{itemize}

\begin{defbox}[Definition --- distinguished opens]
$D(f)=\{\fp\mid f\notin\fp\}=\Spec(R)\mathbin{\backslash} V(f)$, the locus where $f$ is
``nonzero.'' These form a \textbf{basis} for the topology, and crucially
$D(f)\cong\Spec(R_f)$: \emph{localizing the ring $=$ restricting to an
open.}\scite{00E4}{$D(f)$ a basis and $D(f)\cong\Spec(R_f)$; also $\Spec(R)$ is always quasi-compact, Stacks \Tag{00E8}.}
\end{defbox}

\section{Two kinds of point: closed vs.\ generic}

The Zariski topology is \textbf{not Hausdorff} --- points need not be closed,
and that's a feature.

\begin{center}
\begin{tabular}{@{}L{0.24\linewidth} L{0.40\linewidth} L{0.26\linewidth}@{}}
\toprule
\geo{\textbf{Point type}} & \alg{\textbf{Which prime}} &
\textbf{Closure $\overline{\{\fp\}}=V(\fp)$} \\
\midrule
\geo{closed point} & \alg{maximal ideal $\fm$} & just itself \\
\geo{generic point} & \alg{minimal prime (e.g.\ $(0)$ in a domain)} & a whole
irreducible component \\
\bottomrule
\end{tabular}
\end{center}

For a domain, $(0)$ is prime and its closure is the entire space: a single
\term{generic point} whose ``smear'' is everywhere.\scite{00ES}{Irreducible closed sets $\leftrightarrow$ their unique generic points.}
Irreducible closed sets correspond bijectively to their unique generic points
--- geometry and algebra locked together.

\section{See it: four spectra}

\begin{center}
\begin{tabular}{@{}L{0.22\linewidth} L{0.50\linewidth} L{0.20\linewidth}@{}}
\toprule
$R$ & $\Spec(R)$ & \textbf{Shape} \\
\midrule
a field $k$ & one point $(0)$ & a dot \\
$\ZZ$ & generic $(0)$ $+$ one closed point $(p)$ per prime $p$ & ``arithmetic
line'' \\
$\CC[x]$ & generic $(0)$ $+$ closed $(x-a)$ for each $a\in\CC$ & the affine line
$\AA^1$ \\
$k[x]/(x^2)$ & one point $(x)$ --- but with nilpotents & a ``fat'' point \\
\bottomrule
\end{tabular}
\end{center}

\medskip

\noindent The ``arithmetic line'' $\Spec(\ZZ)$ and the affine line
$\Spec(\CC[x])$ share the same shape: a row of closed points, plus one generic
point whose closure is everything.

\begin{center}
\begin{tikzpicture}[>=Stealth,scale=1]
  % base line
  \draw[rulec,thick] (-0.4,0) -- (10.4,0);
  % closed points
  \foreach \x/\lbl in {0.6/(2),1.8/(3),3.0/(5),4.2/(7),5.4/(11),6.6/(13)} {
    \fill[accent] (\x,0) circle (2.2pt);
    \node[below=2pt,font=\footnotesize,accent] at (\x,0) {$\lbl$};
  }
  \node[below=2pt,font=\footnotesize,inksoft] at (7.6,0) {$\cdots$};
  % generic point: drawn as a smeared / open mark to the right, dense everywhere
  \node[above=10pt,font=\footnotesize,accentB] at (9.4,0) {generic point $(0)$};
  \draw[accentB,thick] (8.6,0.0) .. controls (9.0,0.42) and (9.2,0.0) .. (9.4,0.18);
  \fill[accentB] (9.4,0.18) circle (1.6pt);
  \node[above=8pt,font=\footnotesize\itshape,accent] at (3.4,0) {closed points: the primes $(p)$};
  \node[below=20pt,font=\footnotesize,inksoft] at (5.0,0)
       {$\Spec(\ZZ)$: closed points $(2),(3),(5),\dots$ and a dense generic point $(0)$};
\end{tikzpicture}
\end{center}

The last row is the payoff of using \emph{all} primes and remembering the ring:
$k[x]/(x^2)$ and $k$ have the same single point but are different schemes --- the
nilpotent $x$ records an infinitesimal ``thickening'' invisible to the set of
points alone. Hold onto that; it's where schemes beat classical varieties.

\begin{center}
\begin{tikzpicture}[>=Stealth]
  % plain point: Spec k
  \fill[ink] (0,0) circle (2.4pt);
  \node[below=4pt,font=\footnotesize] at (0,0) {$\Spec(k)$: a dot};
  % fat point: Spec k[x]/(x^2)
  \begin{scope}[xshift=5cm]
    \shade[inner color=algc!45,outer color=paper2] (0,0) circle (10pt);
    \draw[algc,thick] (0,0) circle (10pt);
    \fill[algc] (0,0) circle (2.4pt);
    \node[below=14pt,font=\footnotesize] at (0,0)
         {$\Spec\bigl(k[x]/(x^2)\bigr)$: a ``fat'' point};
  \end{scope}
\end{tikzpicture}
\end{center}

\section*{Retrieve it}

Close your notes first --- recall, don't reread. Each answer gives feedback
instantly.

\begin{quizlist}[$\Spec(R)$ --- recall check]
  \quizitem{The points of $\Spec(R)$ are the \underline{\hspace{3em}} of $R$.}
    {prime ideals}{maximal ideals}{radical ideals}{principal ideals}{A}
    {By definition a point is a prime ideal. Maximal ideals give only the closed
     points; $(0)$ in a domain is a point but not maximal.}
  \quizitem{A closed set $V(I)$ depends on the ideal $I$ only through its \underline{\hspace{3em}}.}
    {radical}{nilradical}{generators}{dimension}{A}
    {$V(I)=V(\sqrt I)$, because $\fp\supseteq I \iff \fp\supseteq\sqrt I$.
     Radical ideals biject with closed sets.}
  \quizitem{$\Spec$ of a field has exactly \underline{\hspace{3em}}.}
    {one point}{two points}{no points}{many points}{A}
    {A field's only prime ideal is $(0)$, so its spectrum is a single point.}
  \quizitem{The distinguished opens $D(f)$ form a \underline{\hspace{3em}} for the topology.}
    {basis}{subbasis}{filter}{cover}{A}
    {Every open is a union of $D(f)$'s, so they are a basis --- and
     $D(f)\cong\Spec(R_f)$.}
  \quizitem{If $I\subseteq J$, then $V(I)$ \underline{\hspace{3em}} $V(J)$.}
    {contains}{equals}{misses}{splits}{A}
    {More functions forced to vanish means fewer points: the correspondence
     reverses inclusions, so $V(I)\supseteq V(J)$.}
  \quizitem{For a domain $R$, the generic point of $\Spec(R)$ is \underline{\hspace{3em}}.}
    {the zero ideal}{a maximal ideal}{the unit ideal}{every closed point}{A}
    {In a domain $(0)$ is prime and its closure is the whole space, so $(0)$ is
     the dense generic point.}
  \quizitem{The closed points of $\Spec(\ZZ)$ correspond to \underline{\hspace{3em}}.}
    {the prime numbers}{all the integers}{the zero ideal}{the unit ideals}{A}
    {Maximal ideals of $\ZZ$ are $(p)$ for $p$ prime; $(0)$ is the (non-closed)
     generic point.}
  \quizitem{$\Spec\bigl(k[x]/(x^2)\bigr)$ has exactly \underline{\hspace{3em}}.}
    {one point}{two points}{no points}{four points}{A}
    {The only prime is $(x)$ --- one point --- but the nilpotent $x$ makes it a
     ``fat'' point, a different scheme from $\Spec(k)$.}
\end{quizlist}

\begin{primarybox}
\textbf{\textsc{\textcolor{accentB}{Read next (primary source)}}}\par\smallskip
Ravi Vakil, \emph{The Rising Sea}, the chapter \textbf{``The spectrum of a
ring''} (\S3.2--3.5). Do at least the exercises on computing $\Spec$ for small
rings --- that's where the picture sticks. Ground truth in the Stacks Project,
\Tag{00DZ}.
\end{primarybox}
